\documentclass{../notatki}

\title{Komputery kwantowe \\ fizyka operacji na transmonach}

\begin{document}

\section{Złącze Josephsona}

Jeśli weźmiemy dwa metale przewodzące, w stanie nadprzewodnictwa i umieścimy
między nimi jakiś izolator, to zacznie zachodzić efekt Josephsona. W ramach
niego, powstałe w wyniku nadprzewodnictwa pary Coopera, czyli pary elektronów,
będa przeskakiwać przez izolator. To spowoduje przepływ prądu, mimo izolatora,
oraz braku przyłożonego napięcia.

\subsection{Pary coopera}

Powstawanie par Coopera jest elementem teorii nadprzewodnictwa BCS. Metale
w skali mikroskopowej mają regularne struktury atomowe, które pozwalają na
swobodny przepływ elektronów. W wyniku \textit{niezwykle skomplikowanych}
wpływów fononów i elektronów, elektrony nawzajem się przyciągają.
Lecz, w wysokich (normalnych) temperaturach, energia cieplna jest w stanie
przełamać to niezwykle słabe oddziaływanie. Oczywiście zatem, niska temperatura
nie łamie tych oddziaływań, pozwalając na powstawanie par Coopera.

Oddziaływanie to, jest niezwykle słabe, ale też zaskakująco daleko-dystansowe.
To oznacza, że wiele par Coopera może występować w jednym miejscu.
Przez to, tworzą one kondensat, gdzie tworzą własną meta-strukturę.

Przez tą własną meta-strukturę, złamanie jednej pary wymaga złamania wszystkich
par, stąd ekscytacje pojedynczych elektronów stają się nie możliwe,
oraz wszystkie pary mają ten sam stan kwantowy (energetyczny).

\subsection{Ładunek}

Ładunek przechodzący przez złącze jest dany wzorem:
$$
I = I_c \sin \varphi
$$
$I_c$ to prąd krytyczny, a $\varphi$ to różnica między funkcjami falowymi
dwóch stron złącza.

Podstawową zależnością, która musi zachodzić, jest:
$$
K_BT \leq \hbar \omega_q \leq 2 \Delta
$$

\end{document}
